\documentclass[11pt]{article}
\usepackage[T1]{fontenc}
\usepackage{lmodern}
\usepackage[margin=1in]{geometry}
\usepackage{amsmath,amssymb,amsthm}
\usepackage{enumitem}
\usepackage{graphicx}
\usepackage[hidelinks]{hyperref}
\emergencystretch=2em

\newtheorem{theorem}{Theorem}
\newtheorem{lemma}[theorem]{Lemma}
\newtheorem{corollary}[theorem]{Corollary}
\newcommand{\D}{\mathbb D}
\newcommand{\K}{\mathbb K}
\newcommand{\SC}{\operatorname{SC}}
\newcommand{\Bcal}{\mathcal B}
\newcommand{\Lcal}{\mathcal L}

\title{Two Isometric Classical Subspaces\\
of Supercyclic Operators Beyond $\ell_1$}
\author{Run \texttt{fa\_banach\_001}, lane 8}
\date{July 22, 2026}

\begin{document}
\maketitle

\begin{abstract}
Alves and Souza ask which Banach spaces other than $\ell_1$ admit an
isometric copy of a classical infinite-dimensional Banach space inside the
set of supercyclic operators, apart from zero. We give two explicit positive
cases. First, on $c_0$ the map
\[
 (\lambda_k)_{k\geq1}\longmapsto\sum_{k\geq1}\lambda_kB^k,
\]
where $B$ is the unweighted backward shift, is an isometry from $\ell_1$
whose nonzero values are all supercyclic. This holds over both real and
complex scalars. Second, on every infinite-dimensional separable complex
Hilbert space, the supercyclic operators together with zero contain a closed
isometric copy of the disk algebra $A(\D)$. The first proof uses an exact
finite-support right inverse and the Supercyclicity Criterion. The second
uses adjoint Hardy multipliers and the Godefroy--Shapiro eigenvector
criterion. These results answer two canonical cases of the source question;
they do not classify all Banach spaces.
\end{abstract}

\section{Source question and scope}

The source is T.~R.~Alves and G.~C.~Souza,
\emph{On the set of supercyclic operators}, arXiv:2403.18678, later
published in the \emph{Bulletin of the London Mathematical Society}
\cite{source}. Its Theorem 1.1 proves spaceability for every separable Banach
space and, more sharply, constructs an isometric copy of $\ell_1$ when the
underlying space is $\ell_1$. Question 3.1 on PDF page~4 then asks:
\begin{quote}
Which spaces $X$, other than $\ell_1$, permit the existence of (isometric)
copies of classical infinite-dimensional Banach spaces, excluding the null
operator, within the set of supercyclic operators in $\Lcal(X)$?
\end{quote}
See Figure~\ref{fig:question}.

\begin{figure}[ht]
\centering
\includegraphics[width=.96\textwidth]{figures/open_problem_crop.png}
\caption{Question 3.1 and its immediate context in arXiv:2403.18678,
PDF p.~4.}
\label{fig:question}
\end{figure}

The wording is classificatory. Our two theorems give complete constructions
for $c_0$ and for separable complex Hilbert spaces, but do not characterize
all $X$. We therefore record them as a substantial partial answer.

\section{Proof intuition}

For $c_0$, use exactly the operator series that gives the source's isometric
$\ell_1$ copy on $\ell_1$. The upper operator-norm estimate is the triangle
inequality. The $c_0$ norm gives the reverse estimate even more directly:
one finite unit vector can align the phases of any finite initial block of
coefficients in a single output coordinate.

Supercyclicity is also transparent on finite supports. If the first nonzero
coefficient is $\lambda_p$, factor the operator as $B^p(\lambda_pI+N)$,
where $N$ strictly lowers supports. Thus $N$ is locally nilpotent and a
finite Neumann series inverts $\lambda_pI+N$ on $c_{00}$. Composing this
inverse with the forward shift produces a right inverse. Forward powers of
the original operator eventually annihilate every finitely supported vector,
which verifies the Supercyclicity Criterion without any growth estimate.

For Hilbert space, work on $H^2(\D)$. Adjoint multiplication by a nonconstant
analytic symbol has reproducing kernels as eigenvectors. After rescaling a
nonzero symbol that vanishes at zero, its eigenvalues occur both inside and
outside the unit circle, with dense spans on both sides. The
Godefroy--Shapiro criterion makes the rescaled operator hypercyclic, hence
the original operator supercyclic. A conjugate reflection of the symbol
corrects the conjugate-linearity of taking adjoints and produces a genuinely
complex-linear isometric embedding of the disk algebra.

\section{An isometric \texorpdfstring{$\ell_1$}{ell-1} copy on
\texorpdfstring{$c_0$}{c-0}}

Let $\K$ be $\mathbb R$ or $\mathbb C$, let
\(c_0=c_0(\mathbb N_0;\K)\), and define the backward and forward shifts by
\[
 B(x_0,x_1,x_2,\ldots)=(x_1,x_2,\ldots),
 \qquad
 F(x_0,x_1,\ldots)=(0,x_0,x_1,\ldots).
\]
Write $c_{00}$ for the finitely supported sequences.

\begin{theorem}[The $c_0$ construction]\label{thm:c0}
The map
\begin{equation}\label{eq:Phi-c0}
 \Phi:\ell_1(\K)\longrightarrow\Lcal(c_0),
 \qquad
 \Phi(\lambda)=T_\lambda:=\sum_{k=1}^{\infty}\lambda_kB^k,
\end{equation}
is a linear isometry. Moreover, $T_\lambda$ is supercyclic whenever
$\lambda\neq0$. Consequently,
\[
 \Phi(\ell_1)\setminus\{0\}\subseteq\SC(c_0),
\]
and $\SC(c_0)\cup\{0\}$ contains a closed isometric copy of $\ell_1$.
\end{theorem}

\begin{proof}
Because $\|B\|=1$, the series in \eqref{eq:Phi-c0} converges in operator
norm and
\begin{equation}\label{eq:upper}
 \|T_\lambda\|\leq\sum_{k\geq1}|\lambda_k|=\|\lambda\|_1.
\end{equation}
For $N\geq1$, choose $x^{(N)}\in c_{00}$ with $\|x^{(N)}\|_\infty\leq1$
and, for $j\geq0$,
\[
 x^{(N)}_{j}=\begin{cases}
  \overline{\lambda_{j-1}}/|\lambda_{j-1}|,&
       2\leq j\leq N+1\text{ and }\lambda_{j-1}\neq0,
       \quad(\K=\mathbb C),\\
  \operatorname{sgn}(\lambda_{j-1}),&
       2\leq j\leq N+1\text{ and }\lambda_{j-1}\neq0,
       \quad(\K=\mathbb R),\\
  0,&\text{otherwise}.
 \end{cases}
\]
Thus $x^{(N)}_{k+1}$ cancels the phase of $\lambda_k$ for
$1\leq k\leq N$. Hence
\[
 \|T_\lambda\|\geq |(T_\lambda x^{(N)})_0|
   =\sum_{k=1}^N|\lambda_k|.
\]
Letting $N\to\infty$ and using \eqref{eq:upper} proves that $\Phi$ is an
isometry.

It remains to prove supercyclicity. Fix $0\neq\lambda\in\ell_1$, and put
\[
 p=\min\{k\geq1:\lambda_k\neq0\},
 \qquad
 A=\sum_{j=0}^{\infty}\lambda_{p+j}B^j=\lambda_pI+N,
\]
where $N=\sum_{j\geq1}\lambda_{p+j}B^j$. Then
\begin{equation}\label{eq:factor}
 T_\lambda=B^pA=AB^p.
\end{equation}
On $c_{00}$ the operator $N$ strictly lowers the largest support index, so it
is locally nilpotent. Therefore
\begin{equation}\label{eq:R}
 R y=\lambda_p^{-1}
   \sum_{r=0}^{\infty}(-\lambda_p^{-1}N)^r y,
 \qquad y\in c_{00},
\end{equation}
is a finite sum for each $y$, maps $c_{00}$ into itself, and satisfies
$AR=RA=I$ there. Define $S:c_{00}\to c_{00}$ by
\[
 S=RF^p.
\]
Since $B^pF^p=I$, equations \eqref{eq:factor}--\eqref{eq:R} give
\begin{equation}\label{eq:rightinverse}
 T_\lambda S=B^pARF^p=I\quad\text{on }c_{00}.
\end{equation}

We now apply the Supercyclicity Criterion in the form recalled by the source
\cite{source}. Take both dense test sets to be $c_{00}$, take $n_j=j$, and
take $S_j=S^j$. For fixed $x,y\in c_{00}$, the operator $A$ does not
increase the largest support index. By \eqref{eq:factor},
\[
 T_\lambda^j x=B^{pj}A^j x=0
\]
for all sufficiently large $j$. Consequently
\[
 \|T_\lambda^j x\|\,\|S^j y\|\longrightarrow0.
\]
Also, \eqref{eq:rightinverse} and associativity give
\[
 T_\lambda^jS^j y=y
\]
for every $j$. These are precisely the two clauses of the Supercyclicity
Criterion, so $T_\lambda$ is supercyclic.
\end{proof}

\section{The Hilbert-space construction}

Let $A(\D)$ denote the disk algebra with the supremum norm, and put
\[
 A_0(\D)=\{f\in A(\D):f(0)=0\}.
\]
For $f\in A(\D)$ define $f^\#(z)=\overline{f(\overline z)}$. On the Hardy
space $H^2=H^2(\D)$, let $M_g$ be multiplication by $g$.

\begin{lemma}[Inside/outside eigenvector criterion]\label{lem:GS}
Let $T$ be a bounded operator on a separable complex Banach space. If the
spans of the eigenvectors associated respectively with eigenvalues of modulus
less than one and greater than one are both dense, then $T$ is hypercyclic.
\end{lemma}

\begin{proof}
This is the Godefroy--Shapiro criterion \cite{GS}. On the span of the
eigenvectors with $|\lambda|<1$, $T^n$ tends pointwise to zero. On the span
of eigenvectors with $|\mu|>1$, define the approximate inverse at time $n$
by multiplying a $\mu$-eigenvector by $\mu^{-n}$. These maps tend pointwise
to zero and are exact right inverses for $T^n$. The Hypercyclicity Criterion
applies.
\end{proof}

\begin{theorem}[The Hilbert-space construction]\label{thm:H}
Let $H$ be an infinite-dimensional separable complex Hilbert space. Then
$\SC(H)\cup\{0\}$ contains a closed complex-linear isometric copy of the
disk algebra $A(\D)$.
\end{theorem}

\begin{proof}
It suffices by unitary equivalence to work on $H^2(\D)$. Define
\begin{equation}\label{eq:Phi-H}
 \Psi:A_0(\D)\longrightarrow\Lcal(H^2),
 \qquad
 \Psi(f)=M_{f^\#}^{*}.
\end{equation}
Both $f\mapsto f^\#$ and taking the adjoint are conjugate-linear, so $\Psi$
is complex-linear. Moreover,
\begin{equation}\label{eq:isom-H}
 \|\Psi(f)\|=\|M_{f^\#}\|=\|f^\#\|_\infty=\|f\|_\infty.
\end{equation}
For completeness, the multiplier-norm equality follows from the upper bound
$\|M_g\|\leq\|g\|_\infty$ and the kernel identity
\[
 M_g^*k_w=\overline{g(w)}k_w,
 \qquad
 k_w(z)=\frac{1}{1-\overline wz},quad w\in\D,
\]
which gives the reverse bound after taking the supremum over $w$.

Fix $0\neq f\in A_0(\D)$. Then $f^\#$ is nonconstant and
$f^\#(0)=0$. Choose $\alpha\neq0$ so large that the analytic function
\[
 \psi=\overline\alpha f^\#
\]
has a value of modulus greater than one in $\D$. The sets
\[
 U_- =\{w\in\D:|\psi(w)|<1\},
 \qquad
 U_+ =\{w\in\D:|\psi(w)|>1\}
\]
are nonempty and open: $U_-$ contains a neighborhood of zero, while $U_+$
is nonempty by the choice of $\alpha$. For every nonempty open
$U\subseteq\D$, the kernels $\{k_w:w\in U\}$ have dense span in $H^2$;
indeed, a function orthogonal to all of them vanishes on $U$ and hence
vanishes identically. Since
\[
 \alpha\Psi(f)=M_\psi^*,
 \qquad
 M_\psi^*k_w=\overline{\psi(w)}k_w,
\]
Lemma~\ref{lem:GS} shows that $\alpha\Psi(f)$ is hypercyclic. Its orbit is
contained in the projective orbit of $\Psi(f)$ because
$(\alpha\Psi(f))^n=\alpha^n\Psi(f)^n$. Thus $\Psi(f)$ is supercyclic.

Finally, multiplication by $z$ maps $A(\D)$ onto $A_0(\D)$ and is an
isometry, because the supremum norm of every disk-algebra function is its
supremum on the unit circle. Combining this isometry with \eqref{eq:Phi-H}
and \eqref{eq:isom-H} proves the result on $H^2$, and unitary conjugation
transfers it to $H$.
\end{proof}

\begin{corollary}[Partial answer to Question 3.1]
Question 3.1 has explicit positive answers for $X=c_0$ and for every
infinite-dimensional separable complex Hilbert space. The corresponding
classical spaces can be taken to be $\ell_1$ and $A(\D)$, respectively.
\end{corollary}

\section{Verification, limitations, and novelty}

\paragraph{Verification.}
In Theorem~\ref{thm:c0}, the lower norm estimate uses only finite vectors,
so no forbidden constant-modulus element of $c_0$ is introduced. The inverse
series in \eqref{eq:R} is finite on each input; no convergence of that series
as a bounded operator is asserted or needed. The Supercyclicity Criterion
does not require a uniform estimate on $S^j y$, because $T_\lambda^j x$ is
eventually exactly zero.

For Theorem~\ref{thm:H}, the reflection $f^\#$ is essential: without it,
$f\mapsto M_f^*$ would be conjugate-linear. The two open eigenvalue regions
are nonempty by the vanishing at zero and the scalar rescaling. Density of
their kernel spans follows directly from analytic uniqueness. These checks
are expanded in the accompanying \texttt{VERIFICATION.md}.

\paragraph{Limitations.}
The $c_0$ theorem works over both scalar fields. The disk-algebra theorem is
stated over the complex field, where analytic multiplier and hypercyclicity
arguments apply. No claim is made for nonseparable Hilbert spaces, which
cannot support a supercyclic operator, or for a classification of all
separable Banach spaces.

\paragraph{Novelty check.}
Cheap run/corpus indexes and bounded searches by arXiv id, title, authors,
the exact question, and combinations of ``isometric copy'', $c_0$, Hilbert
space, disk algebra, and supercyclic operators found no later answer. The
only indexed paper citing the source is a 2026 article on a different
frequent-supercyclicity criterion \cite{later}; it does not identify an
isometric classical subspace of $\Lcal(c_0)$ or $\Lcal(H)$. This is a bounded
search, not an exhaustive novelty guarantee. Independent mathematical and
literature review is recommended before external use.

\begin{thebibliography}{9}
\bibitem{source}
T.~R. Alves and G.~C. Souza,
\emph{On the set of supercyclic operators},
Bull. Lond. Math. Soc. \textbf{56} (2024), 2483--2494;
arXiv:2403.18678,
\url{https://arxiv.org/abs/2403.18678}.

\bibitem{GS}
G.~Godefroy and J.~H. Shapiro,
\emph{Operators with dense, invariant, cyclic vector manifolds},
J. Funct. Anal. \textbf{98} (1991), 229--269,
\url{https://doi.org/10.1016/0022-1236(91)90078-J}.

\bibitem{later}
T.~R. Alves, G.~Botelho, and V.~V. F\'{a}varo,
\emph{On frequently supercyclic operators and an
$(\mathcal F,\Gamma)$-supercyclicity criterion with applications},
J. Math. Anal. Appl. \textbf{556} (2026), 130158,
\url{https://doi.org/10.1016/j.jmaa.2025.130158}.
\end{thebibliography}

\end{document}
